\documentclass[11pt,a4paper]{article}
\usepackage{webclases-lesson}
\lessonSetup{T2}{v0.3}{T2 - Movimiento en el plano}
\begin{document}
\shorthandoff{<>}
\color{Ink}
\coverHeader{T2 --- Movimiento en el plano}{Componentes, proyectiles y giros: un mismo reloj, distintos modelos.}
En dos dimensiones estudiamos cada componente y despu\'es reconstruimos el
movimiento. Las herramientas son las de T0 y T1; ahora tambi\'en cambia la direcci\'on.
\begin{lessonindex}
\indexentry{sec:b1}{1. Componer movimientos}
\indexentry{sec:b2}{2. Lanzar y seguir un proyectil}
\indexentry{sec:b3}{3. Entender un giro}
\end{lessonindex}
Usaremos $g=9,81\mpss$ en los tiros, sin resistencia del aire y con gravedad uniforme.
\precisionnote

\block{1}{Componer movimientos}{sec:b1}
La posici\'on $\vec r=x\ii+y\jj$ une el origen con el objeto.
Entre dos instantes, $\Delta\vec r=\vec r_f-\vec r_i$ une los extremos del recorrido;
su m\'odulo no coincide en general con la longitud de la trayectoria.
\[
\vec v=\frac{d\vec r}{dt}=v_x\ii+v_y\jj,\qquad
\vec a=\frac{d\vec v}{dt}=a_x\ii+a_y\jj.
\]
Se deriva e integra por componentes. La velocidad es tangente a la trayectoria,
pero $\vec r$ no tiene por qu\'e serlo. La aceleraci\'on indica el cambio de
velocidad y puede apuntar en otra direcci\'on.

\heading{Qu\'e significa sumar velocidades}
La velocidad de un nadador respecto a la orilla resulta de sumar la velocidad
respecto al agua y la del agua respecto a la orilla:
\[
\vec v_{\mathrm{nadador/orilla}}
=\vec v_{\mathrm{nadador/agua}}+\vec v_{\mathrm{agua/orilla}}.
\]
No se suman siempre sus m\'odulos. Primero se especifican referencias y ejes.

\begin{examplebox}{nadador y corriente}
\origin{Adaptaci\'on del caso del original T2 con los datos de la versi\'on actual.}
\step{Planteamiento} Un r\'io de $42,0\units{m}$ de ancho lleva corriente de
$0,600\mps$. El nadador avanza perpendicular a la orilla a $1,40\mps$ respecto
al agua. Halla tiempo de cruce, deriva y rapidez respecto a la orilla.
\solutionblock
\step{Modelo} Vista en planta: $+x$ sigue la corriente y $+y$ atraviesa el r\'io.
Aqu\'i $y$ no es altura. Las dos componentes comparten tiempo.
\[
\vec v_{\mathrm{orilla}}=(0,600\ii+1,40\jj)\mps,\quad
x=0,600t,\quad y=1,40t.
\]
\begin{center}
\begin{tikzpicture}[x=.085cm,y=.085cm,font=\small]
\fill[BlueSoft] (-4,0) rectangle (48,42);
\draw[BrandBlue] (-4,0)--(48,0) (-4,42)--(48,42);
\draw[axis] (0,0)--(33,0) node[below right] {$x$ (m)};
\draw[axis] (0,0)--(0,49) node[above] {$y$ (m)};
\draw[vector] (0,0)--(18,42);
\draw[guide] (18,0)--(18,42)--(0,42);
\node[left] at (0,42) {$42$};
\node[below] at (18,0) {$18$};
\draw[vectoramber] (28,23)--(42,23) node[midway,above] {corriente};
\end{tikzpicture}
\end{center}
\step{C\'alculo} $t=42/1,40=30,0\units{s}$,
$x=0,600(30)=18,0\units{m}$ y
$|\vec v|=\sqrt{0,600^2+1,40^2}=1,52\mps$.
\step{Interpretaci\'on} Dividir $42$ entre $1,52$ ser\'ia usar una longitud
transversal con una rapidez oblicua. La corriente no acelera el cruce transversal
en este modelo, pero s\'i cambia el punto de llegada.
\variation{Si se duplica solo la corriente, el cruce sigue durando
$30,0\units{s}$ y la deriva es $36,0\units{m}$.}
\end{examplebox}

\block{2}{Lanzar y seguir un proyectil}{sec:b2}
Elegimos $+x$ horizontal, $+y$ hacia arriba y $t=0$ al soltar el objeto.
El modelo de part\'icula sin aire tiene $\vec a=(0,-g)$:
\par\medskip
\begin{tabularx}{\linewidth}{@{}Y Y@{}}
\toprule
Horizontal: MRU & Vertical: MRUA\\
\midrule
$x=x_0+v_{0x}t$ & $y=y_0+v_{0y}t-\tfrac12gt^2$\\
$v_x=v_{0x}$ & $v_y=v_{0y}-gt$\\
$a_x=0$ & $a_y=-g$\\
\bottomrule
\end{tabularx}
\par\medskip
Si la rapidez de salida es $v_0$ y el \'angulo es $\theta$ desde $+x$,
$v_{0x}=v_0\cos\theta$ y $v_{0y}=v_0\sin\theta$.
El suceso pedido decide la ecuaci\'on: para llegar al suelo imponemos $y=0$;
para el m\'aximo, $v_y=0$. Al conocer $t$, calculamos las otras componentes.

\begin{examplebox}{lanzamiento horizontal}
\origin{Adaptaci\'on del tiro horizontal del original T2 con datos de la versi\'on actual.}
\step{Planteamiento} Una bola sale horizontalmente de una mesa de
$1,25\units{m}$ a $3,00\mps$. Halla tiempo de ca\'ida, alcance y velocidad de impacto.
\solutionblock
\step{Modelo} $v_{0y}=0$, no $a_y=0$. Antes de calcular, una bola soltada desde
la misma altura debe tardar lo mismo: ambas comparten su movimiento vertical.
\step{C\'alculo}
\[
0=1,25-4,905t^2\Rightarrow t=0,505\units{s},\qquad
x=3t=1,51\units{m}.
\]
$\vec v_{\mathrm{impacto}}=(3,00\ii-4,95\jj)\mps$ y
$|\vec v|=\sqrt{33,5}=5,79\mps$; forma $58,8^\circ$ bajo la horizontal.
\begin{center}
\begin{tikzpicture}[x=2.2cm,y=2.2cm,font=\small]
\draw[axis] (0,0)--(2.2,0) node[right] {$x$ (m)};
\draw[axis] (0,0)--(0,1.65) node[above] {$y$ (m)};
\draw[BrandBlue,thick,domain=0:1.5145,samples=60] plot(\x,{1.25-4.905*(\x/3)^2});
\foreach \t in {0,.25,.5} {
 \fill ({3*\t},{1.25-4.905*\t*\t}) circle (1.6pt);
 \draw[guide] (0,{1.25-4.905*\t*\t})--({3*\t},{1.25-4.905*\t*\t});
}
\draw[vector] (1.5145,0)--(1.8745,-.594) node[right] {$\vec v$};
\node[left] at (0,1.25) {$1,25$}; \node[below left] at (1.5145,0) {$1,51$};
\end{tikzpicture}
\end{center}
Trayectoria $y(x)=1,25-4,905(x/3)^2$. Los puntos corresponden a $0$, $0,25$ y
$0,50\units{s}$; la flecha de velocidad usa la proporci\'on $3:-4,95$.
\step{Interpretaci\'on} Doblar $v_{0x}$ duplica el alcance a $3,03\units{m}$
y mantiene $t=0,505\units{s}$; cambia la velocidad total de llegada.
\end{examplebox}

\begin{examplebox}{salto oblicuo y condiciones del alcance}
\origin{Adaptaci\'on del modelo de salto del original T2 con datos de la versi\'on actual.}
\step{Planteamiento} Un bal\'on sale del suelo a $12,0\mps$ con
$\theta=35^\circ$ y regresa al mismo nivel. Halla tiempo, alcance y altura m\'axima.
\solutionblock
\step{Modelo} $v_{0x}=12\cos35^\circ$, $v_{0y}=12\sin35^\circ$.
La ra\'iz $t=0$ de $y(t)=0$ es la salida; buscamos la otra ra\'iz.
\step{C\'alculo}
\[
v_{0x}=9,83\mps,\quad v_{0y}=6,88\mps,\quad
t_f=\frac{2v_{0y}}g=1,40\units{s}.
\]
Con componentes sin redondear:
$R_h=v_{0x}t_f=13,8\units{m}$ y
$\Delta h_{\max}=v_{0y}^2/(2g)=2,41\units{m}$.
\step{Interpretaci\'on} En el m\'aximo queda $\vec v=v_{0x}\ii$:
solo se anula la componente vertical. $R_h$ es el alcance, no el radio circular $R$.
\end{examplebox}

Cuando salida y llegada est\'an a la misma altura y $v_{0y}>0$,
\[
t_f=\frac{2v_0\sin\theta}g,\qquad
R_h=\frac{v_0^2\sin2\theta}g.
\]
Son consecuencias del modelo, no f\'ormulas universales para cualquier desnivel.

\heading{Mismo alcance, vuelos distintos}
Con $v_0=20,0\mps$ desde el suelo, $30^\circ$ y $60^\circ$ tienen
$R_h=35,3\units{m}$. A $30^\circ$ el vuelo dura $2,04\units{s}$ y alcanza
$5,10\units{m}$; a $60^\circ$, $3,53\units{s}$ y $15,3\units{m}$.
\begin{center}
\begin{tikzpicture}[x=.18cm,y=.18cm,font=\small]
\draw[axis] (0,0)--(40,0) node[right] {$x$ (m)};
\draw[axis] (0,0)--(0,19) node[above] {$y$ (m)};
\draw[BrandBlue,thick,domain=0:35.311943,samples=80]
 plot(\x,{tan(30)*\x-9.81*\x*\x/(2*400*cos(30)^2)});
\draw[GreenLine!70!black,dashed,thick,domain=0:35.311943,samples=80]
 plot(\x,{tan(60)*\x-9.81*\x*\x/(2*400*cos(60)^2)});
\foreach \h in {5.096840,15.290520} \draw[guide] (0,\h)--(17.655971,\h) node[midway,above] {};
\node[left] at (0,5.1) {$5,10$}; \node[left] at (0,15.29) {$15,3$};
\node[right] at (22,14.3) {$60^\circ$};
\node[right] at (23,4.5) {$30^\circ$};
\node[below] at (35.311943,0) {$35,3$};
\end{tikzpicture}
\end{center}
Las curvas comparten escala espacial y punto de llegada, no instante de llegada.
La complementariedad deja de dar este resultado si cambia la altura final.

\begin{examplebox}{llegar a otra altura y elegir ra\'ices}
\origin{Adaptaci\'on del lanzamiento desde una ventana del original T2 con datos actuales.}
\step{Planteamiento} Una pelota sale desde $8,00\units{m}$ a $10,0\mps$
y $20^\circ$. Calcula impacto en el suelo y los pasos por $y=8,40\units{m}$.
\solutionblock
\step{Modelo} $x=(10\cos20^\circ)t$, $y=8+(10\sin20^\circ)t-4,905t^2$.
No usamos el tiempo de vuelo de alturas iguales.
\step{C\'alculo} Para $y=0$:
\[
t=\frac{10\sin20^\circ\pm\sqrt{(10\sin20^\circ)^2+156,96}}{9,81}.
\]
Las ra\'ices son aproximadamente $1,67$ y $-0,975\units{s}$:
el impacto f\'isico es a $1,67\units{s}$ y $x=15,7\units{m}$.
La altura m\'axima es $8,60\units{m}$, as\'i que s\'i alcanza $8,40\units{m}$.
Al imponer esa altura se obtienen
$t_1=0,149\units{s}$ y $t_2=0,549\units{s}$.
\step{Interpretaci\'on} Pasa por esa altura subiendo y bajando,
en $x_1=1,40\units{m}$ y $x_2=5,16\units{m}$.
Son dos puntos distintos, no dos choques con el mismo balc\'on.
Si el origen vertical se coloca en la ventana, el suelo est\'a en $y'=-8$;
el tiempo de impacto permanece igual.
\end{examplebox}

\heading{Trayectoria espacial frente a gr\'aficas temporales}
Para el bal\'on de $12\mps$ y $35^\circ$, estas gr\'aficas utilizan las
mismas componentes y el mismo intervalo $0\le t\le1,40324\units{s}$.
\begin{center}
\begin{tikzpicture}[font=\small]
\pgfmathsetmacro{\vx}{12*cos(35)}
\pgfmathsetmacro{\vy}{12*sin(35)}
\pgfmathsetmacro{\tf}{2*\vy/9.81}
\begin{scope}[x=3cm,y=.22cm]
\draw[axis] (0,0)--(1.6,0) node[right] {$t$ (s)};
\draw[axis] (0,0)--(0,16) node[above] {$x$ (m)};
\draw[BrandBlue,thick,domain=0:\tf] plot(\x,{\vx*\x});
\node[left] at (0,\vx) {$9,83$}; \node[below] at (1,0) {$1$};
\draw[guide] (1,0)--(1,\vx)--(0,\vx);
\end{scope}
\begin{scope}[xshift=7.4cm,x=3cm,y=1cm]
\draw[axis] (0,0)--(1.6,0) node[right] {$t$ (s)};
\draw[axis] (0,0)--(0,3.4) node[above] {$y$ (m)};
\draw[BrandBlue,thick,domain=0:\tf,samples=60] plot(\x,{\vy*\x-4.905*\x*\x});
\node[left] at (0,2) {$2$}; \node[below] at (1,0) {$1$};
\end{scope}
\begin{scope}[yshift=-5.3cm,x=3cm,y=.22cm]
\draw[axis] (0,0)--(1.6,0) node[right] {$t$ (s)};
\draw[axis] (0,0)--(0,14) node[above] {$v_x$ (m/s)};
\draw[BrandBlue,thick] (0,\vx)--(\tf,\vx);
\node[left] at (0,\vx) {$9,83$}; \node[below] at (1,0) {$1$};
\end{scope}
\begin{scope}[xshift=7.4cm,yshift=-5.3cm,x=3cm,y=.22cm]
\draw[axis] (0,0)--(1.6,0) node[right] {$t$ (s)};
\draw[axis] (0,-8)--(0,10) node[above] {$v_y$ (m/s)};
\draw[BrandBlue,thick,domain=0:\tf] plot(\x,{\vy-9.81*\x});
\node[left] at (0,\vy) {$6,88$}; \node[left] at (0,-\vy) {$-6,88$};
\node[below left] at (.7016,0) {$0,702$}; \node[above] at (1.4032,0) {$1,40$};
\end{scope}
\end{tikzpicture}
\end{center}
La trayectoria $y(x)$ tambi\'en es parab\'olica, pero su eje horizontal son metros.
En las cuatro gr\'aficas anteriores son segundos: $x(t)$ es recta, $y(t)$ es par\'abola,
$v_x(t)$ es constante y $v_y(t)$ tiene pendiente $-g$.

\block{3}{Entender un giro}{sec:b3}
\heading{Radi\'an, periodo y frecuencia}
Un radi\'an es el \'angulo que abarca un arco de longitud igual al radio.
Para un \'angulo con signo $\theta$, $s_{\mathrm{orientado}}=R\theta$.
Una vuelta completa tiene $2\pi$ radianes. Si no cambia el sentido,
la longitud recorrida es $R|\Delta\theta|$; el desplazamiento tras una vuelta es cero.
\[
\omega=\frac{d\theta}{dt},\qquad
\alpha=\frac{d\omega}{dt},\qquad
v_{\theta}=R\omega,\qquad |\vec v|=R|\omega|.
\]
$v_\theta$ es una componente con signo; la rapidez no puede ser negativa.
En MCU, $\omega$ es constante: $\theta=\theta_0+\omega\Delta t$.
El periodo $T$ es el tiempo de una vuelta y la frecuencia $f$ cuenta vueltas por segundo:
$T=1/f$ y $|\omega|=2\pi f$.

\heading{Cambiar de direcci\'on tambi\'en es acelerar}
En una circunferencia, la velocidad es tangente y la aceleraci\'on normal apunta
al centro. La componente tangencial cambia la rapidez:
\[
a_n=\frac{|\vec v|^2}{R}=R\omega^2,\qquad
a_\theta=R\alpha,\qquad
|\vec a|=\sqrt{(R\alpha)^2+(R\omega^2)^2}.
\]
$a_n$ es un m\'odulo; $a_\theta$ sigue el signo angular elegido.
En MCU, $a_\theta=0$ pero $a_n\ne0$. En MCUA, ni el m\'odulo ni la direcci\'on de
la aceleraci\'on total tienen por qu\'e permanecer constantes.
\begin{center}
\begin{tikzpicture}[x=1cm,y=1cm,font=\small]
\draw[BrandBlue] (0,0) circle (1.5);
\foreach \ang in {0,90} {
 \coordinate (P) at (\ang:1.5);
 \fill (P) circle (2pt);
 \draw[vector] (P)--++({\ang+90}:1.1);
 \draw[vectorgreen] (P)--++({\ang+180}:.9);
}
\node[right] at (1.5,1) {$\vec v$};
\node[below] at (.8,0) {$\vec a_n$};
\node[align=left,right] at (3,0) {MCU antihorario: flechas de velocidad\\de igual longitud y distinta direcci\'on.\\Esquema sin escala num\'erica.};
\end{tikzpicture}
\end{center}

\begin{examplebox}{misma rueda, distinto radio}
\origin{Elaboraci\'on propia sobre la comparaci\'on de radios del original T2.}
\step{Planteamiento} Una rueda de radio $0,350\units{m}$ gira a $240\units{rpm}$.
Halla $f$, $T$, $\omega$, rapidez y aceleraci\'on normal del borde; compara con medio radio.
\solutionblock
\step{Modelo} Elegimos como positivo el sentido de giro. Todos sus puntos dan las
mismas vueltas por segundo, aunque recorren distintas longitudes de arco.
\step{C\'alculo}
\[
f=\frac{240}{60}=4,00\units{Hz},\quad T=0,250\units{s},\quad
\omega=2\pi f=25,1\units{rad/s}.
\]
En el borde: $|\vec v|=0,350(8\pi)=8,80\mps$,
$a_n=0,350(8\pi)^2=221\mpss$.
A $R/2$: $|\vec v|=4,40\mps$ y $a_n=111\mpss$.
\step{Interpretaci\'on} A igual $\omega$, ambas magnitudes aumentan con $R$.
A igual rapidez en dos trayectorias diferentes, $a_n=v^2/R$ es mayor en la de menor radio.
Hay que indicar qu\'e se mantiene fijo.
\variation{Duplicar las rpm duplica la rapidez y cuadruplica $a_n$;
no modifica el radio de una rueda r\'igida.}
\end{examplebox}

\begin{examplebox}{frenado angular}
\origin{Adaptaci\'on del MCUA del original T2 con los datos de la versi\'on actual.}
\step{Planteamiento} Un disco pasa de $20,0\units{rad/s}$ al reposo en
$5,00\units{s}$ con $\alpha$ constante. Halla vueltas y aceleraciones a
$R=0,200\units{m}$, inicialmente y a los $2,00\units{s}$.
\solutionblock
\step{Modelo} Radio fijo y sentido positivo igual al giro inicial:
\[
\omega=\omega_0+\alpha t,\qquad
\Delta\theta=\omega_0t+\tfrac12\alpha t^2.
\]
\step{C\'alculo}
$\alpha=(0-20)/5=-4,00\units{rad/s^2}$.
$\Delta\theta=50,0\units{rad}=7,96$ vueltas.
Siempre en este tramo $a_\theta=-0,800\mpss$.
En $t=0$, $a_n=80,0\mpss$; en $t=2$, $\omega=12,0\units{rad/s}$
y $a_n=28,8\mpss$, de modo que $|\vec a|=28,8\mpss$ a tres cifras.
\step{Interpretaci\'on} La aceleraci\'on angular es constante, pero $a_n$ disminuye
al frenar. Si queda parado, el modelo termina en $5,00\units{s}$;
no seguimos hacia $\omega<0$ sin una causa que invierta el giro.
\end{examplebox}

\begin{exercisebox}{elegir una condici\'on}
Con $v_0=20,0\mps$, compara tiros a $20^\circ$, $45^\circ$ y $70^\circ$
desde el suelo hasta el mismo nivel, sin aire. Ordena alcance, altura y tiempo
y justifica en dos frases. Despu\'es indica la direcci\'on de salida de una bolita
que abandona una gu\'ia circular.
\solutionblock
$45^\circ$ tiene el mayor alcance; $20^\circ$ y $70^\circ$ empatan.
Altura y tiempo aumentan en el orden $20^\circ<45^\circ<70^\circ$.
En el instante de abandonar la gu\'ia, la velocidad es tangente, no radial.
La trayectoria posterior depende de las aceleraciones que sigan actuando.
\end{exercisebox}

\heading{Fuentes y procedencia}
El original T2 y su background aportan composici\'on, tiros, alturas distintas y
MCU/MCUA. Se conservan los datos actuales salvo correcciones aritm\'eticas.
Redacci\'on y diagramas propios; comparaci\'on de \'angulos y variaciones de elaboraci\'on propia.
\source{https://fisquiweb.es/Apuntes/apun1B.htm}{FisQuiWeb / IES La Magdalena: inspiraci\'on principal, 1.\textordmasculine\ Bachillerato.}
Contraste externo sobre independencia de componentes y l\'imites de la f\'ormula del alcance:
\source{https://openstax.org/books/university-physics-volume-1/pages/4-3-projectile-motion}{OpenStax, University Physics 1, 4.3: Projectile Motion.}
\end{document}
